\chapter{Coefficient rings and complete local structure}
\label{ch:coefficient-rings}
\renewcommand{\chaptermark}[1]{\markboth{\thechapter.\ #1}{}}

This appendix proves both the regular-ring quotient used for local
duality and the finite regular subring used for nonvanishing. It follows
the coefficient-ring approach in \cite[\S160]{st-alg}. We use
formal lifting and the split conormal criterion from
\cite[FL-01, FL-02, FL-07 and KD-17]{core}, and the completion
properties in \cite[ALG-05-C04-R02--R04 and S01]{p103}.

\section{Lifting fields and a flat lifting criterion}

\begin{lemma}[Two formally smooth lifting facts]\label{res:C01}
A field is formally smooth over its prime field. Also, if
\(A\to B\) is flat, \(K\subset A\) satisfies \(K^2=0\), and
\(B/KB\) is formally smooth over \(A/K\), then \(B\) is formally
smooth over \(A\).
\end{lemma}
\begin{proof}
First let \(F\subset L\) be a finite-type field extension with \(F\)
the prime field. In characteristic zero choose a transcendence basis
and the remaining finite algebraic extension is separable. In
characteristic \(p\), write \(r=\operatorname{trdeg}_F L\).
For any transcendence basis \(t\) the degrees over \(F(t)^p\)
and the Frobenius isomorphism give
\[
 [L:L^p]=[F(t):F(t)^p]=p^r,
\]
since \(F^p=F\) and \([L^p:F(t)^p]=[L:F(t)]\). Adjoin successively
\(x_i\notin L^p(x_1,\ldots,x_{i-1})\). Each step has degree \(p\) because
its minimal polynomial divides \(T^p-x_i^p\), whose monic factors of
degree \(0<e<p\) have coefficient \(-e x_i\) at \(T^{e-1}\).
After \(r\) steps the monomials with exponents \(<p\) form a basis
over \(L^p\). If a nonzero polynomial relation in the \(x_i\) over
\(F\) had minimal degree, grouping its monomials by exponents modulo
\(p\) and taking \(p\)-th roots of the coefficients would yield
\(\sum_\alpha x^\alpha g_\alpha(x)^p=0\).
The basis forces each \(g_\alpha(x)=0\), contradicting minimality for
one nonzero \(g_\alpha\) (a nonzero constant cannot vanish).
These \(x_i\) form a transcendence basis. Write \(F'=F(x_i)\).
Then \(L=L^pF'\). In the finite algebraic extension \(L/F'\), some
\(p^e\)-th power of every element of \(L\) lies in its maximal separable
subfield. Take sufficiently large \(e\) for finitely many algebraic
generators, writing each irreducible polynomial as a separable
polynomial in a \(p\)-power of its variable. But \(L=L^{p^e}F'\) by
iteration, so the extension is separable.

For a square-zero surjection \(T\to T/N\) and a prime-field algebra
map \(L\to T/N\), arbitrary lifts of a separable transcendence basis
give a map on its rational function field, since the images of the
nonzero denominators are units modulo \(N\) and hence units in \(T\).
Lift the remaining finite separable generators successively. If
\(h(b)\in N\) for a lift of a root and \(h'(b)\) is a unit,
replace \(b\) by \(b-h'(b)^{-1}h(b)\). Taylor expansion modulo
\(N^2=0\) gives an exact root. This constructs a lift for the
finite-type subfield.

For an arbitrary field \(L\), choose a lift in \(T\) for each value
of the map on \(L\). Its error in respecting addition, multiplication,
and the prime field is a system of linear equations for corrections in
\(N\), now an \(L\)-vector space via the given map. Indeed products of
two corrections vanish, and multiplication by any chosen lift acts
on \(N\) by its residue. A finite set of equations uses only finitely
many elements of \(L\). The preceding argument lifts their generated
subfield, so that finite subsystem has a solution. More explicitly,
the equations prescribe a linear map from the span of their
finite-support coefficient rows in the free vector space on the
unknowns to \(N\). Finite solvability makes this prescription
well-defined. Extend the map to the entire free vector space by a
basis, giving simultaneous corrections. The resulting lift is an
algebra map. This argument verifies the infinite-generation step.

For the second assertion let \(P=A[T_i]\twoheadrightarrow B\) be
a polynomial presentation with kernel \(J\), allowing arbitrarily many
variables. Flatness gives \(J\cap KP=KJ\) by tensoring the presentation
with \(A/K\). Thus the conormal module of the presentation after
reduction is \((J/J^2)/K(J/J^2)\). The corresponding conormal map
is the reduction of \(d:J/J^2\to\bigoplus_i B\,dT_i\). By formal
smoothness it has a left inverse after reduction. Lift the images
of the free basis to obtain \(s:\bigoplus_i B\,dT_i\to J/J^2\).
Then \(sd=1+e\) with image of \(e\) in \(K(J/J^2)\), and \(e^2=0\).
Thus \((1-e)s\) is a left inverse of \(d\). The split conormal
criterion proves formal smoothness of \(B/A\).
\end{proof}

\section{The coefficient ring}

\begin{lemma}[A Cohen ring with prescribed residue field]\label{res:C02}
For every field \(k\) of characteristic \(p>0\) there is a complete
discrete valuation ring \(\Lambda\) with uniformizer \(p\) and residue
field \(k\). For \(n\ge1\) the map
\(\mathbb Z/p^n\to\Lambda/p^n\Lambda\) is formally smooth.
\end{lemma}
\begin{proof}
Begin with the \(p\)-adic integer ring. In its inverse-limit
description a nonzero element has a first nonzero digit and is a power
of \(p\) times an element invertible modulo every \(p^n\). This proves
the discrete valuation assertion at this starting stage.
Well-order generators of \(k/\mathbb F_p\) and lift the corresponding
tower of residue fields. Suppose \(D\) is a discrete valuation ring
with uniformizer \(p\) and residue field \(k_0\). For a transcendental
adjoining, use \(D[T]_{(p)}\) with residue field \(k_0(T)\).
For an algebraic adjoining, lift its monic irreducible polynomial to
\(D[T]\) and take the monic quotient. This quotient is finite free,
and all its maximal ideals lie over \(p\). Its quotient by \(p\)
is a field, so it is already local. In both cases the new ring
is Noetherian local with maximal ideal generated by the
nonzerodivisor \(p\). By Krull intersection each nonzero element
has a finite \(p\)-order, hence is a power of \(p\) times a unit.
Its ideals are principal by taking the least such order. These steps
are injective and take units to units, so the same finite-order
argument at limit stages, in the directed union, proves that the
union is still a discrete valuation ring with uniformizer \(p\).
Its residue field is the union of the prescribed residue fields.
Complete it. The comparison modulo \(p^n\), faithful flatness, and
Noetherianity of completion are
\cite[ALG-05-C04-S01]{p103}. They give the same maximal ideal and
keep \(p\) regular. Consequently the completion is the required
\(\Lambda\).

The discrete valuation ring \(\Lambda\) is torsion-free over
\(\mathbb Z\), hence flat over \(\mathbb Z\) by the ideal criterion, which
tests the injectivity of multiplication by each positive integer.
Base change makes \(\Lambda/p^n\) flat over \(\mathbb Z/p^n\).
At \(n=1\) apply the field case of Result~\ref{res:C01}. Inductively
\((p^{n-1})\subset\mathbb Z/p^n\) is square-zero when \(n\ge2\).
Apply the second part of Result~\ref{res:C01} to lift formal
smoothness from \(n-1\) to \(n\).
\end{proof}

\begin{theorem}[Coefficient ring of a complete local ring]\label{res:C03}
Let \((R,\mathfrak m,k)\) be a complete Noetherian local ring.
In equal characteristic a coefficient field \(k\to R\) induces the
identity on \(k\). If \(k\) has characteristic \(p>0\) there is also
a continuous map \(\Lambda\to R\), with \(\Lambda\) as in
Result~\ref{res:C02}, inducing the identity on the residue field.
Its kernel is \(0\) if \(p\) is not nilpotent in \(R\), and is
\((p^e)\) if \(p^e=0\) with \(e\) least.
\end{theorem}
\begin{proof}
In residue characteristic zero each nonzero integer is a unit in
\(R\), so we have a coefficient map from \(\mathbb Q\).
Lift the identity \(k\to R/\mathfrak m\) successively to
\(R/\mathfrak m^n\), using Result~\ref{res:C01} on the square-zero
successive quotients. In equal characteristic \(p\), the same proof
uses \(\mathbb F_p\).
For general residue characteristic \(p>0\) start with
\(\Lambda/p\to R/\mathfrak m\). At step \(n\) compose the given map
\(\Lambda/p^n\to R/\mathfrak m^n\) with
\(\Lambda/p^{n+1}\to\Lambda/p^n\).
As a map over \(\mathbb Z/p^{n+1}\) it lifts through
\(R/\mathfrak m^{n+1}\twoheadrightarrow R/\mathfrak m^n\) by
Result~\ref{res:C02}. Passing to inverse limits gives the map.
It preserves the topologies because \(p\in\mathfrak m\). Every
proper nonzero ideal of \(\Lambda\) is a power of \(p\), giving the
kernel assertion. A map from a field inducing the residue identity
is injective.
\end{proof}

\section{Two kinds of regular presentation}

\begin{theorem}[Complete regular covers and presentations]\label{res:C04}
Every complete Noetherian local ring \(R\) is a quotient of a
complete regular local ring \(S\) by a local surjection inducing
the identity on its residue field. A complete regular local ring
of equal characteristic is \(k[[X_1,\ldots,X_d]]\) with \(d=\dim R\).
In mixed characteristic with \(p\notin\mathfrak m^2\) it is
\(\Lambda[[X_1,\ldots,X_{d-1}]]\), with \(\Lambda\) of
Result~\ref{res:C02}.
\end{theorem}
\begin{proof}
Use the coefficient field or complete discrete valuation ring \(D\)
of Result~\ref{res:C03} and send finitely many variables to generators
of \(\mathfrak m\). Evaluation of formal series converges in \(R\).
The source \(D[[X_1,\ldots,X_s]]\) is complete with maximal ideal
generated by the variables and, for a valuation coefficient ring,
by \(p\). To see regularity directly, its associated graded ring is
the polynomial ring over \(k\) on these generators. In degree \(n\)
the coefficient of each monomial in the variables records its
appropriate residue digit at total degree \(n\). These descriptions
also prove completeness and separatedness degree by degree.
A finite set of initial forms generating the initial ideal of any
ideal lifts to finitely many elements of the ideal. Successively
subtract their multiples to increase order. The coefficients
converge, showing those lifts generate the ideal. This uses the
Noetherian polynomial graded ring and gives Noetherianity of the
source. The chains of prime ideals obtained by adjoining generators
one by one give dimension at least the number of generators.
The height theorem gives the reverse bound. Thus the source is
regular local.

To prove that the map onto \(R\) is surjective, lift an element modulo
\(\mathfrak m\), then lift each error in
\(\mathfrak m^n/\mathfrak m^{n+1}\) by a homogeneous polynomial
in the chosen generators. Since the generators come from the
maximal ideal of the source, these corrections converge there.
In equal characteristic and regular dimension \(d\), choose only
\(d\) minimal generators. In the stated mixed characteristic case,
include \(p\) in a minimal set and use \(d-1\) variables for the
others, with \(p\) provided by the coefficient map. The same
successive lifting yields surjectivity. Both source and target have
dimension \(d\). In the regular source, a nonzero ideal has quotient
of dimension at most \(d-1\) by prepending the zero prime to any
chain containing it. Hence the kernel is zero in these two cases.
\end{proof}

\begin{theorem}[Finite regular normalization]\label{res:C05}
A complete Noetherian local domain \((B,\mathfrak m,k)\) of
dimension \(d\) contains a complete regular local subring
\((S,\mathfrak n,k)\) of dimension \(d\) such that \(B\) is finite
over \(S\) and \(\sqrt{\mathfrak n B}=\mathfrak m\).
\end{theorem}
\begin{proof}
In equal characteristic send \(d\) variables over the coefficient
field to a system of parameters. In mixed characteristic the
coefficient discrete valuation ring embeds, since \(B\) is a
domain of characteristic zero. Here \(p\in\mathfrak m\) is
nonzero and \(\dim B/pB=d-1\). The upper bound follows by
prepending the zero prime to chains and the lower bound by
lifting a parameter system of \(B/pB\), adjoining \(p\) and applying
the height theorem. Send \(d-1\) variables over \(\Lambda\) to such
lifts. In each case the image \(J\) of the source maximal ideal
generates an \(\mathfrak m\)-primary ideal in \(B\) and the source
has dimension \(d\) by Result~\ref{res:C04}'s calculation.
The \(J\)-adic and \(\mathfrak m\)-adic topologies agree, since their
powers are cofinal as in Result~\ref{res:L01}.

Choose finite lifts of a \(k\)-basis of \(B/J\).
Their \(S\)-span is all of \(B\). First approximate modulo \(J\),
then express each term in an error in \(J^n\) as a product of
\(n\) image generators times an element of \(B\), and lift its
coefficient modulo \(J\). These lifts give coefficients in \(S\)
of order at least \(n\), so the coefficients converge in \(S\) and
the error tends to zero in \(B\). Thus \(B\) is finite over the
image of the source. Integral extensions preserve dimension
(\cite[ALG-05-C01-S04]{p103}), so that image has dimension \(d\).
A nonzero kernel would lower the dimension of the regular source
as in Result~\ref{res:C04}. The map is injective, proving the result.
\end{proof}
